\documentclass[10pt,a4paper]{article}

\usepackage[T1]{fontenc}
\usepackage[utf8]{inputenc}
\usepackage[a4paper,top=2.1cm,bottom=2.1cm,left=2.2cm,right=2.2cm]{geometry}
\usepackage{amsmath,amssymb}
\usepackage{booktabs,array,tabularx}
\usepackage{enumitem}
\usepackage{xcolor}
\usepackage{titlesec}
\usepackage{fancyhdr}
\usepackage{parskip}
\usepackage[most]{tcolorbox}

\definecolor{ink}{RGB}{28,32,38}
\definecolor{accent}{RGB}{20,72,168}
\definecolor{cmt}{RGB}{110,120,132}
\definecolor{rule}{RGB}{210,216,224}
\definecolor{cardbg}{RGB}{248,250,253}
\definecolor{good}{RGB}{40,125,75}
\definecolor{bad}{RGB}{185,95,20}

\usepackage[colorlinks=true,linkcolor=accent,urlcolor=accent]{hyperref}

\titleformat{\section}{\Large\bfseries\color{ink}}{\thesection.}{0.5em}{}
\titleformat{\subsection}{\large\bfseries\color{accent}}{}{0em}{}
\titlespacing*{\section}{0pt}{1.0em}{0.4em}
\titlespacing*{\subsection}{0pt}{0.8em}{0.3em}

\pagestyle{fancy}\fancyhf{}
\renewcommand{\headrulewidth}{0.4pt}
\renewcommand{\headrule}{\hbox to\headwidth{\color{rule}\leaders\hrule height \headrulewidth\hfill}}
\fancyhead[L]{\small\color{cmt}GARF — Reading the output · cheat sheet}
\fancyhead[R]{\small\color{cmt}\thepage}

\newcommand{\code}[1]{\texttt{#1}}
\newcommand{\opt}[1]{\texttt{\textbf{#1}}}
\newcommand{\good}{\textcolor{good}{\textbf{Good}}}
\newcommand{\bad}{\textcolor{bad}{\textbf{Concern}}}
\newcommand{\act}{\textbf{Action}}

% ---- one card per object ---------------------------------------------------
% \begin{card}{Name}{stands for}{where it appears}
%   \what{...} \howread{...} \interp{...}
% \end{card}
\newtcolorbox{card}[3]{enhanced,breakable,colback=cardbg,colframe=accent!60,boxrule=0.5pt,arc=2pt,
  left=6pt,right=6pt,top=4pt,bottom=2pt,before skip=5pt,after skip=5pt,
  before upper={\parskip=0pt\raggedright},
  title={\textbf{#1}\quad{\normalfont\small #2}\hfill{\normalfont\scriptsize #3}},
  coltitle=white,colbacktitle=accent!85,fonttitle=\small}
\newlength{\cardlab}\setlength{\cardlab}{2.35cm}
\newcommand{\cardrow}[2]{\par\noindent\hangindent=\cardlab\hangafter=1%
  \makebox[\cardlab][l]{\small\bfseries\color{accent}#1}{\small #2}\par\vspace{2pt}}
\newcommand{\what}[1]{\cardrow{What}{#1}}
\newcommand{\howread}[1]{\cardrow{How to read}{#1}}
\newcommand{\interp}[1]{\cardrow{Interpret}{#1}}

\begin{document}

% ---------------- title ----------------
\thispagestyle{empty}
\begin{center}
{\LARGE\bfseries\color{ink} GARF — Reading the output}\\[0.35em]
{\large\color{accent} A cheat sheet: every quantity, what it means, how to read it, how to interpret it}\\[0.6em]
{\normalsize \texttt{garf} (matching) \;·\; \texttt{DIAGNOSTICS} (design checks) \;·\; \texttt{garfte} (estimation)}\\[0.3em]
{\color{cmt}\small GARF-binary · version 1.2 · companion to the \emph{Installation \& Use Guide}}\\[0.5em]
{\small Federico Fantechi, PhD · University of Palermo · \texttt{federico.fantechi@unipa.it}}
\end{center}
\vspace{0.2em}
\hrule height 0.4pt
\vspace{0.5em}

\noindent\textbf{How the sheet works.} One card per quantity. The card's title bar gives its name,
what the acronym stands for and where it appears; then three rows: \textbf{What} it is, \textbf{How to
read} it (scale, range, what a value means), and how to \textbf{Interpret} it (\good, \bad, \act).
The three parts follow the order of work: \textbf{1}~match with \code{garf}; \textbf{2}~check the design
with \code{DIAGNOSTICS}; \textbf{3}~estimate with \code{garfte}. Section~\ref{sec:cheat} puts the
rules of thumb on one page. Examples (``e.g.'') come from a firm panel: 311{,}398 firms, 2009--2019,
$T_0=2016$, 423 treated, $k=25$, caliper 0.2.

\noindent\textbf{Golden rule.} Parts 1 and 2 use pre-treatment information only: revise the design
(covariates, $k$, caliper, \opt{wpower()}, replacement) as often as needed \emph{there}. Once you look
at Part 3, the design is fixed --- changing it because of the estimated effect is specification search.

\noindent\textbf{Notation.} Treated unit $i$, donor (matched control) $j$, period $t$, onset $T_0$,
matching window of $L$ periods before $T_0$, block = a treated unit and its donors, $w_{ij}$ donor
weights ($\sum_j w_{ij}=1$ in each block).

\vspace{-0.4em}
\tableofcontents

% =====================================================================
\section{\texttt{garf}: the matching}\label{sec:garf}

\subsection{Sample and design matrix}

\begin{card}{Units dropped}{incomplete pre-window}{engine log}
\what{Units removed because the outcome or a time-varying covariate is missing in at least one period
of the matching window $[T_0-L,\,T_0-1]$. Only units complete on \emph{all} of them enter the match.}
\howread{A count, printed as a warning. Compare it with the units received.}
\interp{\good: a small share. \bad: a large share (e.g.\ 217{,}719 of 311{,}398, 70\%), or treated units
dropped more than controls --- the matched population is ``units fully observed on these variables''.
\act: find the variables with most gaps; drop them from \opt{tvarying()} or shorten \opt{npre()}.}
\end{card}

\begin{card}{Features}{pre-treatment design matrix}{engine log · \code{features.csv}}
\what{What the forest sees. Per time-varying variable (outcome included): the $L$ lags and 7 trajectory
descriptors --- \code{mean}, \code{delta} (last $-$ first), \code{growth} (symmetric, in $[-2,2]$),
\code{slope}, \code{vol} (s.d.\ of changes), \code{curv} (late minus early slope), \code{dip} (last minus
mean of the others); plus one column per time-invariant covariate.}
\howread{``units $\times$ features''. With $L=5$: 12 features per time-varying variable (e.g.\ $12\times12+3=147$).}
\interp{Matching is on levels \emph{and} dynamics. The outcome is just the first time-varying variable
(no need to list it again). Trees use thresholds: heavy-tailed amounts are better as ratios, logs or asinh.}
\end{card}

\begin{card}{Treated / Control / Treated matched}{block counts}{summary · \code{r(n\_treated)}, \code{r(n\_treated\_matched)}}
\what{Units that entered the match, and treated units with at least one donor.}
\howread{Matched $<$ treated only with a caliper: treated units with no control above it get no donor.}
\interp{The effect estimated later is the effect on the \emph{matched} treated. \bad: many unmatched
(e.g.\ 106 of 423), or unmatched systematically different (\code{blocks.csv}, \code{matched}=0). \act:
report it; compare matched vs unmatched on pre-treatment characteristics; consider a looser caliper.}
\end{card}

\subsection{Similarity and weights}

\begin{card}{Proximity}{forest (leaf) proximity}{\code{matches.csv} · \code{proximity\_minmax}}
\what{Similarity of a donor to a treated unit as seen by the forest: how often (weighted by leaf size)
the two fall in the same leaf. \code{proximity\_minmax} rescales it to $[0,1]$ over the run.}
\howread{0 = least similar pair in the run, 1 = most similar pair. Relative to the run, not comparable across runs.}
\interp{The basis of donor ranking, the caliper and the weights. Many donors near 0: poor overlap (card \emph{Overlap}).}
\end{card}

\begin{card}{Co-occurrence}{share of trees with a shared leaf}{\code{cooccurrence}}
\what{Share of trees in which donor and treated unit end in the same leaf.}
\howread{Between 0 and 1; not rescaled, so comparable across designs that share the same forest.}
\interp{Higher = closer in every direction the forest cares about. Input of FMQ.}
\end{card}

\begin{card}{Caliper}{minimum normalised proximity}{option \opt{caliper()}}
\what{Donors with \code{proximity\_minmax} below the caliper are not allowed.}
\howread{On the normalised $[0,1]$ scale; 0 = no caliper.}
\interp{Trades bias (closer donors) for coverage (unmatched treated). Choose it on the frontier and on
\emph{FMQ all treated}, not on matched-only statistics.}
\end{card}

\begin{card}{weight}{design weight, $w_{ij}\propto\text{proximity}^\gamma$}{variable \code{weight} · option \opt{wpower()}}
\what{Treated rows: 1. Donor rows: proximity to the power $\gamma$, normalised to sum to 1 in the block.}
\howread{$\gamma=0$ equal weights; $\gamma=1$ (default) proportional to proximity; $\gamma>1$ concentrates on the closest.}
\interp{Summing to 1 makes the counterfactual a weighted \emph{mean}, cancels shocks common to all units,
and gives every treated unit the same mass in the ATT. Use as \code{[pw=weight]}. Non-negativity: no
extrapolation outside the donors' range.}
\end{card}

\subsection{Quality of the match}

\begin{card}{SMD}{Standardised Mean Difference}{summary · \code{r(mean\_smd\_matched)} · \code{balance.csv} · 02, 04}
\what{For each feature: treated mean minus comparison mean, divided by the raw pooled s.d. (the same
scale before and after matching). \emph{Raw}: vs all controls; \emph{matched}: vs weighted donors.}
\howread{$|\text{SMD}|$ in s.d.\ units. Summary: mean and max over all features.}
\interp{\good: $<0.10$. \bad: $0.10$--$0.25$ (needs adjustment), $>0.25$ (a weighted comparison will
not remove it). Read the max, not only the mean (e.g.\ max 0.174). \act: 02 tells which variable and
whether level or dynamics.}
\end{card}

\begin{card}{ESS}{Effective Sample Size (of donors)}{summary · \code{r(mean\_ess)} · \code{blocks.csv} · 07}
\what{$\text{ESS}_i=1/\sum_j w_{ij}^2$: how many equally weighted donors the weights are worth.}
\howread{From 1 (one donor carries everything) to the number of donors (equal weights). E.g.\ 5.3 of 25.}
\interp{The \emph{variance} side of the design. ESS well below $k$ is normal with proximity weights.
\bad: 1--2 --- a thin, noisy counterfactual. \act: lower \opt{wpower()} or raise $k$.}
\end{card}

\begin{card}{FMQ}{Forest Match Quality}{summary · \code{r(fmq\_mean)}, \code{r(fmq\_mean\_all)} · 09, 12}
\what{$\text{FMQ}_i=(S_i-\bar c_i)/(S_i^{\max}-\bar c_i)$: weighted co-occurrence with the donors ($S_i$),
relative to random controls ($\bar c_i$) and to the best control in the data ($S_i^{\max}$).}
\howread{0 = no closer than random controls; 1 = every donor as close as the best available. \emph{p10}: the
10\% worst blocks. \emph{All treated}: unmatched count as 0.}
\interp{The \emph{bias} side of the design, and \emph{relative}: compare designs sharing the same forest,
never an absolute target. Rises mechanically with fewer, closer donors --- read it with ESS (frontier).
With a caliper compare on \emph{all treated} (e.g.\ 0.662 matched vs 0.496 all).}
\end{card}

\begin{card}{Twin FMQ}{FMQ scaled by the best twin (earlier definition)}{log · \code{blocks.csv} · 09}
\what{Earlier version of FMQ, kept for comparison: co-occurrence relative to the closest control only.}
\howread{Small numbers on large pools; not comparable with FMQ.}
\interp{Secondary. Use FMQ.}
\end{card}

\begin{card}{Distinct donors / reuse}{size of the comparison group}{summary · 07}
\what{How many different controls serve as donors; with replacement one control can serve several treated.}
\howread{Compare with matched $\times\,k$: e.g.\ 890 donors for $317\times25=7{,}925$ slots $\approx$ 9 uses each.}
\interp{Reuse is not wrong (each treated keeps its best donors), but estimates rest on fewer controls,
blocks become dependent (the \code{gap} SEs correct for it) and \code{method(did)} needs a match without replacement.}
\end{card}

\subsection{How separable the groups are}

\begin{card}{AUC}{Area Under the ROC Curve (out-of-bag)}{summary · \code{r(oob\_auc)} · 10}
\what{Probability that the forest ranks a random treated unit above a random control, on units it did not train on.}
\howread{0.5 = indistinguishable; 1 = perfectly separable.}
\interp{0.5--0.7: little selection, good overlap. 0.7--0.85: real selection, matching matters.
\bad: $>0.85$--$0.9$ (e.g.\ 0.907): treated easy to pick out, overlap limited --- expect unmatched
treated under a caliper and some low-FMQ blocks; check Overlap. Describes the data, not a GARF failure.}
\end{card}

\begin{card}{OOB bal.acc / acc}{Out-Of-Bag balanced accuracy / accuracy}{summary · \code{r(oob\_bacc)}, \code{r(oob\_score)}}
\what{Hit rates out of bag: balanced = mean of the hit rates in the two groups (threshold 0.5); accuracy = overall share right.}
\howread{Balanced accuracy stays near 0.5 when treatment is rare; accuracy approaches the control share.}
\interp{Read the AUC first. Accuracy 0.995 with 99.5\% controls says nothing.}
\end{card}

% =====================================================================
\section{\texttt{DIAGNOSTICS}: checking the design}\label{sec:diag}

Run after \code{garf}, with the matched panel in memory: \code{do "\$GARF\_HOME/DIAGNOSTICS/00\_master.do"}.
Check the settings it prints first. Numbers refer to the do-files; tables go to the log, figures to
\code{./garf\_diagnostics}.

\subsection{Balance}

\begin{card}{Balance by component}{SMD by variable $\times$ levels/dynamics}{02 · log table}
\what{Per variable: mean and max $|$SMD$|$ raw and matched, separately for \emph{levels} (lags, mean) and
\emph{dynamics} (delta, growth, slope, vol, curv, dip); count above 0.10, 0.25, and of features that got worse.}
\howread{One row per variable and component.}
\interp{Locates the imbalance the summary only counts. Levels balanced but dynamics not: same height,
different motion --- a warning for \code{gap} and for parallel trends. Many features worse on one variable:
the forest barely uses it (see 10).}
\end{card}

\begin{card}{VR}{Variance Ratio}{03 · log table}
\what{Variance among treated over the window divided by the variance in the comparison group, raw and matched.}
\howread{1 = same spread.}
\interp{\good: 0.5--2. \bad: outside --- donors do not cover the treated's range, typically the tails,
even when means are balanced.}
\end{card}

\begin{card}{SMD over time}{balance by period}{04 · \code{fig\_balance\_time.pdf}}
\what{SMD by period and variable: raw (solid) and matched (dashed), with $\pm0.1$ and the window start marked.}
\howread{Inside the window: matched by construction. \emph{Before} the window: out of sample.}
\interp{\good: dashed line near 0 also before the window --- the match captures the history, not only $L$
periods of it. \bad: dashed line drifting away the further back --- treated and donors converged only recently.}
\end{card}

\begin{card}{Outcome path}{treated vs counterfactual vs all controls}{05 · \code{fig\_outcome\_path.pdf}}
\what{Mean outcome up to the end of the window: treated (solid), GARF counterfactual (dashed), all controls (grey).}
\howread{Grey--solid distance = selection; dashed--solid distance = what is left after matching.}
\interp{\good: close in the window, parallel before. A constant dashed--solid gap is a level difference:
harmless for \code{did}, fatal for \code{gap}.}
\end{card}

\subsection{The placebo}

\begin{card}{Placebo gap}{design backdated by $H$ periods}{06 · \code{fig\_placebo.pdf}}
\what{Your design rerun with the onset moved back by $H$ (default 3, lowered automatically if the data are
too short), on pre-treatment data. True effect in the pseudo-post period: zero.}
\howread{Treated minus counterfactual (dashed) vs treated minus all controls (grey) across window and pseudo-post.}
\interp{The most direct test of the design. \good: dashed line flat into the pseudo-post period. \bad: a
gap that opens --- the same would happen after the real $T_0$ and look like an effect. E.g.\ data from 2009,
$T_0=2016$, $L=5$: $H=2$, onset 2014.}
\end{card}

\subsection{Donors and overlap}

\begin{card}{$\omega_j$, top shares}{implied donor weight and concentration}{07 · log tables}
\what{$\omega_j=\sum_i w_{ij}$: total weight of donor $j$ across blocks (they sum to the number of treated).
Pool summaries: effective number of donors $(\sum\omega)^2/\sum\omega^2$; share of weight on the top 1, 2, 10.}
\howread{Reuse = number of blocks a donor serves; $\omega$ = how much it counts overall.}
\interp{\bad: a few donors carrying a large share (top 10 above about 20--30\%) --- every estimate depends
on them. \act: look at who they are; expect wider \code{gap} intervals.}
\end{card}

\begin{card}{Overlap}{best proximity · overlap mass}{08 · \code{fig\_overlap.pdf} · \code{blocks.csv}}
\what{Per treated unit: \code{best\_proximity\_minmax} (normalised proximity of its best donor) and
\code{overlap\_mass} (share of trees with at least one control in its leaf).}
\howread{Histograms over treated units; the log lists the weakest and the unmatched.}
\interp{\bad: mass near 0 --- treated units with no close control; a caliper drops them, otherwise their
effect is extrapolated. With caliper $c$, units left of $c$ on the first histogram are the unmatched.}
\end{card}

\begin{card}{FMQ components}{$S_i$, $\bar c_i$, $S^{\max}_i$}{09 · log table}
\what{\code{match\_cooccurrence} ($S_i$), \code{random\_cooccurrence} ($\bar c_i$), \code{best\_cooccurrence}
($S^{\max}_i$) behind each FMQ; the five lowest-FMQ blocks are listed.}
\howread{The three numbers explain a low FMQ.}
\interp{High $\bar c_i$: a \emph{typical} unit, any control is close --- not a problem. Low $\bar c_i$, low $S_i$,
low overlap mass: an \emph{isolated} unit --- weak counterfactual.}
\end{card}

\subsection{The forest and the choice of design}

\begin{card}{MDI importance}{Mean Decrease in Impurity}{10 · log}
\what{The fifteen features the forest relies on most to tell treated from controls, plus the out-of-bag fit.}
\howread{Group them by variable and by type (level vs dynamic).}
\interp{What selection into treatment runs on. Dynamics on top: selection on trajectories --- where GARF has
its edge over matching on levels. One variable dominating: needs a substantive explanation.}
\end{card}

\begin{card}{Sensitivity grid}{19 designs + yours}{11 · \code{sens.dta}}
\what{\code{garf} rerun varying one choice at a time: $k$, $\gamma$, caliper, replacement, seed; plus your design.
Per design: SMD, ESS, FMQ, donors, matched. Skipped on pools of 20{,}000+ units unless \code{GARF\_GRID on}.}
\howread{Compare rows; seeds 1--5 show forest noise.}
\interp{\good: stable balance across seeds and small changes. \bad: balance collapsing for a small change ---
fragile design. Choose the design here, before any estimate.}
\end{card}

\begin{card}{Frontier}{FMQ vs ESS across designs}{12 · \code{fig\_frontier.pdf}}
\what{FMQ (all treated) against mean ESS (log scale), one point per design; your design is a filled square.}
\howread{The $k$ line is the trade-off: moving right buys effective sample size with closeness.}
\interp{\good: your design on or above the line --- little FMQ given up for a large gain in ESS.
Designs above the line beat the $k$ trade-off.}
\end{card}

% =====================================================================
\section{\texttt{garfte}: the estimation}\label{sec:te}

\begin{card}{Gap}{treated minus own counterfactual}{the building block of both methods}
\what{$G_{it}=Y_{it}-\sum_j\tilde w_{ij}Y_{jt}$, with $\tilde w$ the weights renormalised over the donors observed in $t$.}
\howread{In outcome units; one per treated unit and period.}
\interp{Everything \code{garfte} reports is an average of $G$ or of $G$ minus its window mean.}
\end{card}

\begin{card}{\code{method(did)}}{matching Difference-in-Differences}{\code{garfte, method(did)}}
\what{$\tau_i(e)=G_{i,T_0+e}-\bar G_i$, $\bar G_i$ = mean gap over the matching window. Synthetic-DiD analogue.}
\howread{Effect net of any constant level difference left by the match.}
\interp{Assumes parallel trends conditional on the match. Needs a match \emph{without} replacement
(checked). The default for several outcomes and for heavy-tailed variables.}
\end{card}

\begin{card}{\code{method(gap)}}{the gap in levels}{\code{garfte, method(gap)}}
\what{$\tau_i(e)=G_{i,T_0+e}$. Synthetic-control analogue.}
\howread{Effect in levels; window rows show the leftover level difference.}
\interp{Assumes no unmatched level or trend difference. Credible only if the window gap is small (next cards).
Replacement allowed.}
\end{card}

\begin{card}{ATT, ATT($e$)}{Average Treatment effect on the Treated (by event time)}{table · \code{e(b)} · \code{e(att)}}
\what{ATT($e$): simple average over treated units of $\tau_i(e)$. ATT: equal-weight average of the post ATT($e$).}
\howread{In outcome units, per treated unit (each counts once, whatever its number of donors).}
\interp{The shape of the post ATT($e$) (jump, build-up, fade) often says more than the overall ATT. Weighting treated
units by size is a different estimand: average \code{tau} from the units file with those weights.}
\end{card}

\begin{card}{$e$, period}{event time and period label}{table}
\what{$e$ = year $-\,T_0$ ($e=0$: first treated period). Labels: \code{pre} (before the window, never matched on),
\code{window} (matched on), \code{pre (window)} (window, outcome not matched on), \code{post} (effects).}
\howread{\code{pre(\#)}, \code{post(\#)} trim the range shown.}
\interp{\textbf{Window rows are not a pre-trends test} (\code{did}: they sum to 0; \code{gap}: small by construction).
\textbf{Pre rows are}: near 0 with \code{did}, near the window level with \code{gap}.}
\end{card}

\begin{card}{Window gap}{mean gap in the matching window}{\code{gap} footer · \code{e(mean\_window\_gap)} · \code{pre\_gap\_mean}}
\what{Average pre-treatment level difference between treated and their donors.}
\howread{In outcome units, against the overall ATT.}
\interp{\good: small relative to the ATT. \bad: a sizable share --- part of the ``effect'' was there before
treatment; read \code{did}. E.g.\ window gap 53, pre gap 51, post 104--257: a stable level difference of about 50.}
\end{card}

\begin{card}{SE, CI, $p$, df}{Standard Error, Confidence Interval, $t$ degrees of freedom}{table · \code{e(V)} · \opt{level()}}
\what{Conditional on the matched design. Clustered by block; \code{gap} adds the correction for donors shared
between blocks. Student $t$ with (treated units $-1$) df; \opt{level()} default 95.}
\howread{As usual; the CI is also the band in the graph.}
\interp{\bad: fewer than about 30 treated units (warning) --- cluster-type SEs can be too small. Many outcomes:
adjust for multiple testing (Holm, Romano--Wolf) or name a primary outcome.}
\end{card}

\begin{card}{N}{treated units with a defined gap at $e$}{table}
\what{How many treated units contribute to ATT($e$).}
\howread{Falls when units exit or have missing outcomes.}
\interp{A falling N changes the population behind each ATT($e$). \bad: if treatment affects survival.
\act: compare with a balanced subset.}
\end{card}

\begin{card}{Renormalised / skipped cells}{missing donor-periods}{notes · \code{n\_donors\_obs}}
\what{Cells where some donors are missing (weights renormalised over those observed) or where no gap exists (skipped).}
\howread{Counts in the notes; per cell in the units file.}
\interp{Fine when few cells and few donors are involved; check \code{n\_donors\_obs} where many are.}
\end{card}

\begin{card}{tau (units file)}{unit-level effect}{\code{garfte\_<method>\_units.csv}}
\what{One effect per treated unit and period, with \code{pre\_gap\_mean}, \code{n\_donors}, \code{ess}, \code{fmq}.}
\howread{Group or regress on characteristics of the treated (fixed at or before $T_0$).}
\interp{Noisy: their spread overstates true heterogeneity, regression coefficients are fine. With shared donors
they are correlated (plain OLS SEs too small). Check mean vs median: a few units may drive the ATT.}
\end{card}

\begin{card}{Event-study graph}{ATT($e$) with confidence band}{\code{garfte\_<method>.gph/.pdf}}
\what{Line with \opt{level()}\% band, window shaded, dashed line at the onset, overall ATT in the subtitle.}
\howread{Left of the window: pre periods; shaded: window; right of the dashed line: effects.}
\interp{Flat before the window (\code{did}: at 0; \code{gap}: at the window level). A widening band late in the
panel reflects fewer, more dispersed treated units.}
\end{card}

\begin{card}{did vs gap}{bracketing}{run both on a match without replacement}
\what{Under fairly general conditions the two bracket the true effect.}
\howread{Compare their post ATT($e$) and their window/pre rows.}
\interp{Equal, small window gap, flat pre rows: robust. \code{gap} $>$ \code{did} with a stable pre gap: level
difference --- read \code{did}. \code{did} pre rows trending: parallel trends in doubt (check 04, 06). Placebo gap too:
the effect is not credible.}
\end{card}

% =====================================================================
\section{Rules of thumb on one page}\label{sec:cheat}

\renewcommand{\arraystretch}{1.12}
\begin{tabularx}{\textwidth}{@{}>{\raggedright\arraybackslash}p{0.34\textwidth}>{\raggedright\arraybackslash}p{0.2\textwidth}X@{}}
\toprule
Quantity & \good & \bad\\
\midrule
Units dropped (incomplete window) & small share & large share; treated dropped more\\
Treated matched / treated & close to 1 & unmatched differ systematically\\
SMD matched, mean & $<0.10$ & $>0.10$\\
SMD matched, max & $<0.10$--$0.25$ & $>0.25$\\
VR, Variance Ratio (03) & 0.5--2 & outside\\
AUC (out-of-bag) & 0.5--0.8 & $>0.85$--$0.9$: limited overlap\\
OOB balanced accuracy & read only with AUC & ---\\
ESS per treated & several donors & 1--2\\
FMQ & compare designs, with ESS & never on its own\\
SMD before the window (04) & near 0 & drifting\\
Placebo gap (06) & flat & opening in pseudo-post\\
Weight on top 10 donors (07) & small share & a few donors carry the design\\
Overlap (08) & little mass near 0 & many treated with no close control\\
\code{garfte}: window gap (\code{gap}) & small vs ATT & sizable share of ATT\\
\code{garfte}: pre rows (\code{did}) & near 0 & trending\\
\code{garfte}: treated units (clusters) & $\ge 30$ & fewer\\
\code{garfte}: N over $e$ & stable & falling (exits)\\
\bottomrule
\end{tabularx}

\vfill
\begin{center}\small\color{cmt}
GARF-binary · version 1.2 · Reading the output · Federico Fantechi, University of Palermo\\
Rules of thumb, not tests: read every number in the context of the design and the data.
\end{center}

\end{document}
